% v208 locator refresh: re-sync Japanese page/line guide after final English Marked-PDF publisher; translated prose unchanged.

編集者および二名の査読者の詳細かつ建設的な報告に感謝する。指摘を受け、学習アルゴリズムを修正し、未使用の型付け成分を削除し、終端規則に関する完全性証明の不備を補修し、原稿の構成を大幅に改めた。

本回答書で特に断りなく述べる節・定理・命題・補題・系の番号は、いずれも改訂稿の番号である。投稿時の番号には「投稿時」または「旧」と付す。査読者の原文の節構成に対応させた小見出し（例えば「第7節」「第8--10節」）は、投稿時原稿の構成を指す。編集者の意見を E.1、査読者の意見を R1.1--R1.37 および R2.1--R2.11 として番号付けした。引用部分では元の節構成を維持し、体裁の正規化のみを行い、省略箇所は省略記号で示す。関連する意見は各節で主題ごとにまとめた。各回答は改訂後の節または番号付きの定理類を直接示す。以下の参照表には、修正表示PDFにおいて検証したページおよび余白行番号を記す。

\section*{改訂点の概観}
\noindent\textbf{中心的な修正。}
中心的な修正は四つである。第一に、集合依存の再構成作用素 $\Bh$ と保守的な逐次学習器 $\Ah$ を分離し、Gold 学習器の文法表現が構文的に安定するようにした。第二に、未使用の外側文脈の型注釈を削除し、出現文脈で添字付けしていた仮説非終端記号を商によりまとめ、観測された各部分語 $x$ に対し一つの記号 $\NT{x}$ とした。その結果、バッチ構成は Yoshinaka の部分語文法（Clark による紹介もある）に沿い、単項規則の窓比較を $h$-型の一致で置き換えた。第三に、投稿時の特徴標本に暗黙に含まれていた型付き終端規則の証拠語を明示的に列挙し、完全性証明では必要に応じて語彙規則 (L) より先に単項規則 (U) を使うことで、証明の欠落を補修した。第四に、固定 $h$ の代入可能性は非空の内部因子について述べ、$\lambda$ は開始記号で別途扱うこととした。

\noindent\textbf{先行研究と貢献。}
改訂序論では2014年の著者プレプリント（arXiv:1409.6247v3）を研究史上の資料として引用し，
固定 $h$ の定性的な正例学習可能性が Yoshinaka (2015) の一般的な多ソート枠組みからも得られることを認める．
学習器の存在，局所自明半群の古典的判定条件，固定窓シグネチャ自体を新規性とはしない．
本稿の独自の貢献は符号長で測る特徴データ上界（窓・有限群型付けおよび線形の場合を含む），
型付けの変更に関する閉包性，分離結果にある．関連研究として Clark--Yoshinaka (2016)，
Coste--Garet--Nicolas (2012)，Coste--Nicolas (2019)，Costa (2002) を加えた．

これに伴い、主張の適用範囲も絞った。提示順序に依存しないのはバッチ作用素のみであり、逐次学習器ではない。多項式特徴データ上界は特定の定量的条件下でバッチ作用素に対して述べ、逐次学習器については Gold 収束と多項式時間更新を別に証明した。投稿時の第7節の表題 \emph{The Linear Subclass: A Full Polynomial Time-and-Data Theorem} を逐次学習器の特徴づけとして用いることを撤回した。改訂後の定理8.3は、選んだ線形 CFG に対する $\Bh$ の多項式特徴データを示すが、$\Ah$ の逐次的な多項式データ保証は主張しない。第2.1節は、集合ベースの基準（Eyraud--Heinz--Yoshinaka, 2016, Definition 2.8）と、より強い逐次的基準（同 Definition 2.5）との違いを説明する。外側文脈の型が必要だという従来の主張は撤回し、商を取った構成では独立の文脈輸送規則も不要になった。

新たな型付き厚みの分離例を導入した．命題6.3では，
サイズと通常厚みが $O(n)$ の有限 Clark--Eyraud 代入可能言語族 $L_n$
を与える．自明型付けの $\Bh$ には多項式特徴データがあるが，
固定された二元半束型付け $h_{\mathrm U}$ の $\Bh$ では，
長さ $3\cdot2^n-2$ の語が任意の特徴標本に必要となる．
命題6.3(iv)は，同じ型付けクラスの二標的を区別するため，
双方に特徴標本をもつ任意の集合駆動型作用素について
同じ特徴標本下界が成立することを示す．
系6.4は，サイズによらない
$\tau_{h_{\mathrm U}}^{\mathrm{typ}}(G)\ge3\cdot2^n-2$
という任意の縮約済み SSBNF 文法表示の下界を与える．
補題7.6は $U_1\notin\mathbf{LI}\vee\mathbf{Grp}$ を示す．
これと命題6.3を合わせることで，系7.5の代数的仮定を
一様な特徴データ上界から単純に外せないことが分かる．
ただしこの条件の一般的必要性や，任意のGold学習器の実行時間・
収束時刻に対する下界は主張しない．
これらは今回の改訂で追加したものである．ただし補題7.6は，議論を完結させるために記載した初等的な半群論の事実である．

原稿全体も大幅に短縮・再構成した．基礎定義を第2節に集約し、完全性証明用の標的側の精密化よりも前に再構成作用素を定義した。従来の三つの表現力・境界に関する節を第9節に統合し、正規化の技術的証明は付録A--Bに移し，新しい括弧付き言語族の代入可能性の詳細証明を付録Cに置いた。付録Cでは命題6.3(i)を，可逆な決定的トップダウン解析器によって直接証明する．接尾部の保存と文法認識を明示し，完全な部分木の交換可能性と分岐同期から，同一の両端状態を結ぶ経路が構文合同になることを示す．明示的な経路の帰納法を用い，二分節点の子の境界をまたぐ因子も扱う．
投稿時の付録Aにあった共通部分と逆準同型像についての閉包性は、消去される因子を区別する新しい非空因子の規約に合わせて、命題3.2へ簡約した。第6--8節では定量的議論を再構成し、非正規な例と線形部分クラス内の正規フィルターの議論で、線形定理に関する査読者2の疑問に答えた。第9節全体も短縮され、査読者1の提案を受けた合同文脈自由言語クラスに関する命題9.7を加えた。

定理8.3の線形文法の上界は、正規化後の線形文法 $G_0$ について
$\|\mathsf{Wit}(\widetilde G_0)\|=O(|G_0|^3|M|^4)$
という符号長評価として明示できる。ただし、正規化前の入力文法のサイズに対して同じ次数を主張しない。また、非消去逆像に関する観察、分解 $\mathsf{S}^n=\mathsf{S} E(\mathsf{S})\mathsf{S}$、および新設の補題7.1で重要節点の幅の評価を与えた。命題9.1の非線形例では、許容される入口高度集合 $\mathcal Q(x)$ を標準的な補完により定義し、連結に関する恒等式を脚注で裏付け、$ba$ を含む因子と $a^*b^*$ に属する因子を分けて固定 $h_\star$ の代入可能性を証明した。上限付きカウンターの分離に必要な鋭い条件 $\rho\ge2$ について、$\mathrm{CTR}_1$ が $(1,1)$-代入可能であることを脚注で説明した。

改訂で新しく加えた内容、投稿時から存在した結果、Eilenberg (1974)、Kanazawa (1996, 1998)、Yoshinaka (2008)、Pin (1997, 2025) の先行結果は区別した。とくに命題3.5は、局所自明半群の既知の特徴づけを、正語像と核の記法で述べ直すものであり、半群論の新定理ではない。序論には、剛性文法の特徴標本と $k$-値範疇文法の文字列言語に関する有限弾性の議論について、Kanazawa (1998, Chapters 5 and 8) の引用も加えた。これらは相補的な正例学習の結果であって、本論文の固定 $h$ の定理がそこから従うわけではない。予備事項と命題3.1では、正規言語の構文合同と有限モノイドによる特徴づけについて Eilenberg (1974, Chapter III, Propositions 10.1--10.2) を引用した。局所自明半群の結果は引き続き Pin に帰する。

\subsection*{投稿時から改訂後への主要番号対照}
原稿の構成を大幅に変えたため、報告書で頻繁に言及された項目の対応を示す。
\begin{center}\small
\begin{tabular}{@{}p{0.46\linewidth}p{0.46\linewidth}@{}}
\toprule
投稿時の原稿 & 改訂原稿\\\midrule
定義2.3（$\sim_h$）& 定義2.1\\
定義3.2（Clark--Eyraud）& 第3.1節の番号なしの特殊例\\
命題3.5 & 命題3.3\\
定理3.7 & 定理3.4\\
例4.3 & 削除\\
旧規則(2),(3) & 商構成と単項規則(U)\\
補題5.2(i) & 証拠集合(2)、定理5.3\\
定理5.7 & 定理5.4\\
系5.8 & 系5.5\\
命題7.2 & 命題8.1\\
定理7.9 & 定理8.3\\
定理8.6 & 定理9.2（証明簡略化）\\
系8.7 & 削除\\
命題10.3--10.5 & 命題9.1に統合\\
補題9.11 & 補題9.6\\
付録A（閉包性）& 命題3.2（規約変更に対応）\\
$\mathrm{CCL}_{p},\mathrm{CCL}$ & $\mathrm{CTR}_{\rho},\mathrm{CTR}$\\
カウンター記号 $d,u,;$ & $d\mapsto\uparrow$（増加），$u\mapsto\downarrow$（減少），$;\mapsto\#$（リセット）\\
$L_Y$（古典的言語）& $L_{\mathrm{Luk}}$（Łukasiewicz言語）\\
命題9.2, 9.6, 9.9・補題9.8 & 第9.3節の短い記述・引用\\
\bottomrule
\end{tabular}
\end{center}
